% =============================================================================
\section{Timing Design \& Non-Blocking Scheduling}
\label{sec:timing}
% =============================================================================

\subsection{System Timing Specifications}

All temporal dimensions in the washing machine control system are configured in
\code{washing_machine_config.h} and synthesized deterministically from a single $1.0\,\text{kHz}$
hardware time base (Table~\ref{tab:timing}).

\begin{table}[H]
  \centering
  \caption{System timing constants and physical derivations}
  \label{tab:timing}
  \small
  \begin{tabularx}{\linewidth}{lL{2.5cm}Y}
    \toprule
    \textbf{Timing Constant} & \textbf{Nominal Value} & \textbf{Physical Source / Engineering Derivation} \\
    \midrule
    \code{WM_TICK_RATE_HZ}          & $1000\,\text{Hz}$ ($1\,\text{ms}$)  & Derived from Arm Cortex-M7 SysTick ($400\,\text{MHz} / 400{,}000$). \\
    \code{WM_DEBOUNCE_TIME_MS}      & $30\,\text{ms}$    & Symmetric integration window exceeding mechanical chatter ($5\text{--}20\,\text{ms}$). \\
    \code{WM_STOP_WINDOW_MS}        & $1500\,\text{ms}$  & Intentional double-press window under ISO 9241 ($T_{\text{double}} \le 1.5\,\text{s}$). \\
    \code{WM_GESTURE_DOUBLE_GAP_MS} & $250\,\text{ms}$   & Maximum release-to-press interval for K1 double-click token. \\
    \code{WM_GESTURE_HOLD1_MS}      & $600\,\text{ms}$   & K1 Hold Stage 1: $+50\,\text{\textcent}$ deposit in Standby; STOP \#1 in active cycle. \\
    \code{WM_GESTURE_HOLD2_MS}      & $1200\,\text{ms}$  & K1 Hold Stage 2: Confirm in Standby; STOP \#2 in active cycle ($\Delta t = 600\,\text{ms} < 1.5\,\text{s}$). \\
    \code{WM_BLED_RUNNING_HZ}       & $1.0\,\text{Hz}$   & $500\,\text{ms}$ \code{ON} / $500\,\text{ms}$ \code{OFF} ($50\%$ duty square wave). \\
    \code{WM_RLED_ERROR_HZ}         & $2.0\,\text{Hz}$   & $250\,\text{ms}$ \code{ON} / $250\,\text{ms}$ \code{OFF} (high-urgency optical alert). \\
    \code{WM_DEFAULT_CYCLE_TIME_S}  & $1800\,\text{s}$   & Standard 30-minute wash cycle (scaled to 30\,s for demo build). \\
    \code{WM_SPIN_DURATION_S}       & $300\,\text{s}$    & Final $1/6$ of cycle ($1500 < t \le 1800\,\text{s}$): centrifugal spin + drain pump. \\
    \bottomrule
  \end{tabularx}
\end{table}

\subsection{Single Time-Base Paradigm and Zero-Blocking Discipline}

Embedded software in control applications frequently suffers from jitter, responsiveness degradation,
and race conditions when blocking delays (e.g., \code{HAL_Delay()}) or nested interrupt service routines (ISRs)
are introduced. If a $500\,\text{ms}$ blocking sleep were executed to modulate an indicator LED, user button
clicks and coin insertion pulses arriving during that window would be dropped.

To guarantee zero latency and non-blocking determinism, the system architecture decouples hardware time
generation from event processing:
\begin{enumerate}
  \item \textbf{Atomic Interrupt Tick Generation:} The SysTick ISR executes in exactly four CPU instructions:
        \begin{equation}
          \code{void SysTick_Handler(void) \{ g_millis++; \}}
        \end{equation}
        It operates strictly as an atomic counter increment, taking $<15\,\text{ns}$ at $400\,\text{MHz}$.
        No peripheral drivers, FSM transitions, or calculations are executed inside the ISR, eliminating interrupt
        latency and priority inversion.
  \item \textbf{Decoupled Super-Loop Catch-Up:} The main execution loop continuously monitors the difference
        between the hardware timestamp \code{board_millis()} and the software accumulator \code{processed_ms}.
        Every elapsed millisecond is serviced sequentially in a non-blocking catch-up loop (Figure~\ref{fig:timebase}).
\end{enumerate}

\begin{figure}[H]
  \centering
  \resizebox{0.95\linewidth}{!}{%
  \begin{tikzpicture}[
      b/.style={draw, thick, rounded corners, align=center, font=\small, minimum height=0.9cm, fill=white},
      a/.style={-{Stealth[length=2.2mm]}, thick}]
    \node[b, fill=gray!10] (st) at (0,0) {SysTick ISR\\\scriptsize $1.0\,\text{kHz}$ Hardware Tick};
    \node[b] (ms) at (3.8,0) {Atomic Counter\\\scriptsize \code{g_millis}};
    \node[b] (loop) at (8.0,0) {Super-Loop Catch-Up\\\scriptsize \code{while (processed != now)}};
    \node[b] (p1) at (12.4,0) {Discrete Engine\\\scriptsize \code{process_1ms()}};
    \node[b] (stop) at (8,-2.0) {STOP Window\\\scriptsize $w \gets w - 1$};
    \node[b] (bl) at (10.6,-2.0) {Blinker Engine\\\scriptsize $1\,\text{Hz} / 2\,\text{Hz}$ Phase};
    \node[b] (db) at (13.4,-2.0) {Integrator Debounce\\\scriptsize $30\,\text{ms}$ Filter};
    \node[b] (acc) at (16.2,-2.0) {Second Prescaler\\\scriptsize $acc_{\text{ms}} == 1000$?};
    \node[b, fill=stRun] (t) at (16.2,-4.0) {FSM Down-Counter\\\scriptsize $t \gets t - 1$};
    \draw[a] (st) -- (ms);
    \draw[a] (ms) -- (loop);
    \draw[a] (loop) -- (p1);
    \draw[a] (p1) -- (stop.north);
    \draw[a] (p1) -- (bl.north);
    \draw[a] (p1) -- (db.north);
    \draw[a] (p1) -- (acc.north);
    \draw[a] (acc) -- node[right, font=\scriptsize]{\code{TICK_1S}} (t);
  \end{tikzpicture}}
  \caption{Single time-base architecture: Deterministic derivation of all software timers from SysTick}
  \label{fig:timebase}
\end{figure}

\subsection{Mathematical Proof of Zero Cumulative Timing Drift}
\label{sec:drift-proof}

\begin{thmbox}[Zero Cumulative Drift under Non-Preemptive Super-Loop Catch-Up]
  Let $t_{\text{hw}}(k) \in \mathbb{N}$ denote the monotonic millisecond count generated by the hardware
  SysTick ISR. Even in the presence of variable execution delays during graphics rendering
  ($\Delta t_{\text{SPI}} \approx 16.38\,\text{ms}$) or UART packet transmissions, the cumulative timing error
  $\epsilon_{\text{drift}}$ of the 30-minute operational timer relative to the hardware clock satisfies:
  \begin{equation}
    \lim_{t \to \infty} \epsilon_{\text{drift}}(t) = 0.000\,\text{ms}
  \end{equation}
\end{thmbox}

\paragraph{Proof:}
Let $N_{\text{ticks}}(T)$ be the total number of $1\,\text{ms}$ ticks produced by the hardware clock over an
arbitrary wall-clock interval $[0, T]$. Let $M(t)$ be the discrete execution delay of an arbitrary super-loop
iteration (for instance, pushing a $25.6\,\text{KB}$ framebuffer over SPI4 with $M(t) = 16\,\text{ms}$).

During the interval $[t, t + M(t)]$, the hardware counter increments freely:
\begin{equation}
  \Delta g_{\text{millis}} = M(t)
\end{equation}
Upon completion of the blocking operation, the catch-up loop is entered:
\begin{lstlisting}
while (processed_ms != current_ms) {
    processed_ms++;
    process_1ms();
}
\end{lstlisting}
The loop executes exactly $M(t)$ iterations. In each iteration, \code{process_1ms()} advances the millisecond
accumulator:
\begin{equation}
  acc_{\text{ms}} \gets (acc_{\text{ms}} + 1) \pmod{1000}
\end{equation}
emitting the second-level event \code{WM_EVT_TIMER_TICK_1S} whenever $acc_{\text{ms}}$ wraps from $999$ to $0$.

The total number of 1-second countdown ticks dispatched to the FSM over interval $T$ is:
\begin{equation}
  N_{1\text{s}}(T) = \left\lfloor \frac{\sum_{k=1}^K M_k}{1000} \right\rfloor = \left\lfloor \frac{N_{\text{ticks}}(T)}{1000} \right\rfloor
\end{equation}
Because no millisecond tick can ever be skipped, bypassed, or dropped, the total accumulated discrete ticks
match the continuous integral of the hardware clock:
\begin{equation}
  \epsilon_{\text{drift}}(T) = \big| N_{\text{ticks}}(T) - \code{processed_ms}(T) \big| = 0
\end{equation}
Therefore, cumulative drift is mathematically zero. The sole source of real-world timing variance is the
physical quartz crystal stability ($\pm 20\,\text{ppm}$), inducing a worst-case deviation of less than:
\begin{equation}
  |\Delta t_{\text{physical}}| \le 1800\,\text{s} \times 20 \times 10^{-6} = 0.036\,\text{s} \quad (\text{over 30 minutes})
\end{equation}
which is negligible for commercial washing operations. \hfill $\blacksquare$

\subsection{Modulo Phase Preservation in Optical Blinker Drivers}

A common bug in embedded blinker drivers is phase drift or duty-cycle corruption when large discrete
time steps ($\Delta t_{\text{step}} > 1\,\text{ms}$) are processed. In our implementation (\code{hal_led_blinker.c}),
the instantaneous optical state is computed via modular phase arithmetic:
\begin{equation}
  \text{State}_{\text{BLED}}(t) = \left\lfloor \frac{t_{\text{elapsed}}}{500} \right\rfloor \pmod 2
  \label{eq:blink1hz}
\end{equation}
\begin{equation}
  \text{State}_{\text{RLED}}(t) = \left\lfloor \frac{t_{\text{elapsed}}}{250} \right\rfloor \pmod 2
  \label{eq:blink2hz}
\end{equation}
where $0$ corresponds to \code{OFF} and $1$ corresponds to \code{ON}.

Even if a sudden burst of processing delays the super-loop by $28\,\text{ms}$, evaluating
Equations~\eqref{eq:blink1hz} and~\eqref{eq:blink2hz} reconstructs the exact optical state without
accumulating phase error. Changing the operating mode resets $t_{\text{elapsed}} := 0$, guaranteeing
that entering \st{RUNNING} always initiates with an immediate, full $500\,\text{ms}$ illuminated pulse.

\subsection{Double-Press STOP Timing Geometry}

Figure~\ref{fig:stop-window} illustrates the temporal geometry of the dual-strike \code{STOP} detection window.

\begin{figure}[H]
  \centering
  \resizebox{0.95\linewidth}{!}{%
  \begin{tikzpicture}[x=0.55cm, y=0.8cm, font=\scriptsize]
    \draw[-{Stealth}] (0,0) -- (25,0) node[right]{$t$ (s)};
    \foreach \x/\l in {0/0, 5/0.5, 10/1.0, 15/1.5, 20/2.0} { \draw (\x,0.1) -- (\x,-0.1) node[below]{\l}; }
    % case A
    \node[anchor=east] at (-0.3,2.4) {Case A: Valid Dual Strike};
    \draw[fill=stRun!30, draw=stRun, thick] (0,2.0) rectangle (15,2.8);
    \node at (7.5,2.4) {Active Window $w = 1500\,\text{ms}$};
    \draw[thick, -{Stealth[length=2mm]}] (0,3.6) -- (0,2.85) node[pos=0, above]{STOP \#1 ($n \gets 1, w \gets 1500$)};
    \draw[thick, -{Stealth[length=2mm]}] (8,3.6) -- (8,2.85) node[pos=0, above]{STOP \#2 ($t = 0.8\,\text{s} < 1.5\,\text{s}$) $\implies$ Force Abort};

    % case B
    \node[anchor=east] at (-0.3,0.9) {Case B: Expired Window};
    \draw[fill=stPause!30, draw=stPause, thick] (0,0.5) rectangle (15,1.3);
    \node at (7.5,0.9) {Window Expired ($w = 0 \implies n \gets 0$)};
    \draw[thick, -{Stealth[length=2mm]}] (18,1.8) -- (18,0.9) node[pos=0, above right]{STOP at $1.8\,\text{s}$: Treated as New STOP \#1};
  \end{tikzpicture}}
  \caption{Timing geometry of the double-press STOP confirmation window}
  \label{fig:stop-window}
\end{figure}

\why[Switch Debouncing Interval: Why $30\,\text{ms}$?]{%
  Physical pushbutton switches (such as tactile switch K1 on the WeAct PCB) rely on mechanical metal domes.
  Upon closure, elastic kinetic energy causes the contacts to bounce physically for $3\text{--}18\,\text{ms}$,
  generating multiple high-frequency electrical pulses that would register as spurious coin deposits or rapid
  button presses. A software integrator threshold of $30\,\text{ms}$ provides a conservative $1.67\times$
  margin over the maximum observed mechanical chatter duration ($18\,\text{ms}$). Simultaneously, $30\,\text{ms}$
  is well below the human tactile threshold of perceptible sluggishness ($50\text{--}100\,\text{ms}$),
  delivering instantaneous tactile response while completely eliminating false triggers.
}
